\subsection{EABC-Signierung als Richtungs-Symmetrieoperator}
\label{subsec:eabc_signierung_richtungsoperator}

Zur Untersuchung der Rolle der Kleinschen Vierergruppe im quaternionischen Pianouniversum vergleichen wir zwei Versionen des Zustandsmodells:
\begin{enumerate}
	\item eine \emph{unsignierte} Version, in der die quaternionischen Achsenbeiträge allein durch \(\nu_2(N)\), \(\nu_3(N)\), \(\nu_5(N)\) und \(\log(R_{235}(N)+1)\) bestimmt werden,
	\item eine \emph{EABC-signierte} Version, in der dieselben Beiträge durch klassenspezifische Vorzeichenmuster moduliert werden.
\end{enumerate}

Die verwendeten Signaturen sind
\[
e=(+,+,+),\qquad
a=(+,-,-),\qquad
b=(-,+,-),\qquad
c=(-,-,+).
\]

\begin{bemerkung}
	Die EABC-Signierung verändert im vorliegenden Modell nicht die Amplitudenstruktur, sondern primär die räumliche Orientierung der Zustände.
	Sie ist daher als diskreter Richtungs- bzw.\ Polarisationsoperator zu lesen.
\end{bemerkung}

Numerisch ergibt sich für \(N\le 300\), dass die Klasse \(e\) unter der Signierung praktisch invariant bleibt, während die Klassen \(a\), \(b\) und \(c\) deutliche Umorientierungen im Zustandsraum erfahren.
Für die mittlere räumliche Verschiebung zwischen unsignierter und signierter Version erhält man:
\[
\Delta_e \approx 0,\qquad
\Delta_a \approx 9.67,\qquad
\Delta_b \approx 11.17,\qquad
\Delta_c \approx 3.17.
\]

Dies zeigt, dass \(e\) einen fixen Sektor des Modells bildet, während \(a\), \(b\) und \(c\) echte Richtungsmodulationen erzeugen.

Auch in der Projektion auf die \((i,j)\)-Ebene wird die Wirkung der Signierung sichtbar:
\begin{itemize}
	\item \(e\) bleibt im Mittel nahezu unverändert,
	\item \(a\) wird effektiv gespiegelt,
	\item \(b\) wird stark in einen getrennten Winkelbereich verschoben,
	\item \(c\) wird ebenfalls verschoben, jedoch schwächer als \(a\) und \(b\).
\end{itemize}

\begin{proposition}
	Im quaternionischen Pianouniversum wirkt die EABC-Signierung als diskreter Symmetrieoperator der Richtungen.
	Sie trennt die inneren Klassen \(e,a,b,c\) nicht primär durch Stärke, sondern durch Orientierung im Zustandsraum.
\end{proposition}

\begin{proof}[Begründung]
	Die numerischen Vergleiche zeigen, dass sich die mittleren Amplituden der signierten und unsignierten Zustände nicht unterscheiden, während die räumlichen Projektionen in \((i,j,k)\)-Richtung deutlich auseinanderlaufen.
	Damit ist die Wirkung der EABC-Signierung orientierungs- und nicht amplitudengetrieben.
\end{proof}

Von besonderem Interesse ist außerdem die Unterscheidung zwischen \emph{neuen Generatoren} und \emph{harmonischen Nachläufern}. 
Die Vergleichsplots legen nahe, dass beide zwar häufig in derselben groben generatorischen Linie liegen, durch die EABC-Signierung jedoch verschieden stark auseinandergezogen werden.
Dies spricht für eine zusätzliche interne Feinstruktur innerhalb derselben Restgeneratorfamilie.

\begin{figure}[t]
	\centering
	\includegraphicssafe[width=.9\textwidth]{pianouniversum_vergleich_ij_uns_vs_sgn.png}
	\caption{Vergleich der \((i,j)\)-Projektionen des quaternionischen Pianouniversums ohne und mit EABC-Signierung. Die Signierung wirkt vor allem als Umorientierung der Klassen \(a,b,c\), während \(e\) weitgehend invariant bleibt.}
	\label{fig:pianouniverse_ij_compare}
\end{figure}

\begin{figure}[t]
	\centering
	\includegraphicssafe[width=.9\textwidth]{pianouniversum_vergleich_verschiebung.png}
	\caption{Räumliche Verschiebung zwischen unsignierter und signierter Version als Funktion von \(N\). Besonders deutlich betroffen sind die nichttrivialen EABC-Klassen sowie neue Generatoren.}
	\label{fig:pianouniverse_displacement}
\end{figure}

\begin{figure}[t]
	\centering
	\includegraphicssafe[width=.9\textwidth]{pianouniversum_vergleich_generator_nachlaeufer.png}
	\caption{Mittlere Verschiebung von Generatoren und Nachläufern unter EABC-Signierung. Die Vierergruppenstruktur wirkt damit nicht nur familienbildend, sondern auch fein auftrennend innerhalb generatorischer Linien.}
	\label{fig:pianouniverse_generators_followers}
\end{figure}

Insgesamt ergibt sich damit folgendes Bild:
Ohne Signierung erhält man ein schalen- und restgetriebenes Universum, mit Signierung dagegen ein \emph{gruppenpolarisiertes Kugeluniversum}, in dem die Klassen \(e,a,b,c\) als orientierende Sektoren sichtbar werden.